% !PRINTED 302
\weberpage{302}

2. $\mu$ ungerade:

\[
\begin{array}{ll}
\text{a)}\ f^{(\nu-1)}(x),\ f^{(\nu)}(x), & \text{b)}\ f^{(\nu-1)}(x),\ f^{(\nu)}(x) \\[2pt]
\delta_1,\quad + \quad - & \quad - \quad - \\[2pt]
\delta_2,\quad + \quad + & \quad - \quad +
\end{array}
\]

Es geht also bei geradem $\mu$ zwischen $f^{(\nu-1)}$ und $f^{(\nu)}$ kein Zeichenwechsel verloren, bei ungeradem $\mu$ geht ein Zeichenwechsel verloren oder es wird einer gewonnen. Es ist also der Verlust an Zeichenwechseln beim Durchgang durch $\xi$ in der Reihe

\weq{5}{f^{(\nu-1)},\ f^{(\nu)},\ f^{(\nu+1)}\ .\ .\ .\ f^{(\nu+\mu)}}

bei geradem $\mu$ gleich $\mu$, bei ungeradem $\mu$ gleich $\mu \pm 1$, also immer eine \gesperrt{gerade Zahl} und nie negativ.

Wir können diesen Ergebnissen einen übersichtlichen analytischen Ausdruck geben, der ihre Bedeutung besser erkennen lässt.

Wir bezeichnen mit $V(x)$ wie früher die Anzahl der Zeichenwechsel in der Reihe (1) für einen bestimmten Werth von $x$, jedoch mit der näheren Bestimmung, dass, wenn für einen Werth von $x$ einige der Glieder der Reihen verschwinden, diese bei Abzählung der Zeichenwechsel einfach übergangen werden. Wir können dann $V(x)$ als eine Function von $x$ auffassen, die sich aber nur um ganze Zahlen ändern kann, also unstetig ist für die Werthe von $x$, für welche einige Glieder der Reihe (1) verschwinden. Wir wollen dann mit $V(x-0)$ und $V(x+0)$ die Werthe der Function $V(x)$ unmittelbar vor und unmittelbar nach einer solchen Stelle bezeichnen. Dann zeigt die Betrachtung von (4), dass in allen Fällen

\weq{6}{V(\xi + 0) = V(\xi)}

ist, und dass, wenn $\xi$ eine $\mu$-fache Wurzel von $f(x)$ ist,

\weq{7}{V(\xi - 0) = V(\xi) + \mu + 2h,}

worin $h$ eine nicht negative ganze Zahl ist.

Daraus ergiebt sich also, wenn wir zunächst daran festhalten, dass für $x = \alpha$ und $x = \beta$ keine der Functionen (1) verschwindet, dass $V(\alpha) - V(\beta)$ sich von der Summe aller Zahlen $\mu$, d.\,h. von der Anzahl der im Intervall $(\alpha, \beta)$ gelegenen Wurzeln, jede nach dem Grade ihrer Vielfachheit gerechnet, um eine gerade, nicht negative Zahl unterscheidet, und damit ist also das Theorem I von seiner Beschränkung befreit.
